%%%%%%%%%%%%%%%%%%%%% chapter.tex %%%%%%%%%%%%%%%%%%%%%%%%%%%%%%%%%
%
% sample chapter
%
% Use this file as a template for your own input.
%
%%%%%%%%%%%%%%%%%%%%%%%% Springer Nature %%%%%%%%%%%%%%%%%%%%%%%%%%
\motto{Mathematics is, in its own way, the poetry of logical ideas. ---Albert Einstein}
\chapter{Mathematical Foundation}
\label{mathematics}

\abstract*{
  The following chapter will explain the mathematical notation that will 
  be used in the rest of the book. Additionally, it will dive deep into 
  the mathematical formulation of various properties intrinsic to spacetime
  and elaborate the inner workings of integration methods. The reader is 
  required to know basic linear algebra and analysis or otherwise this book 
  would get too complex. For the proof of the reader's understanding exercises
  shall be provided at the end of the chapter, which will be mostly of 
  albegraic nature.
}

\abstract{
  The following chapter will explain the mathematical notation that will 
  be used in the rest of the book. Additionally, it will dive deep into 
  the mathematical formulation of various properties intrinsic to spacetime
  and elaborate the inner workings of integration methods. The reader is 
  required to know basic linear algebra and analysis or otherwise this book 
  would get too complex. For the proof of the reader's understanding exercises
  shall be provided at the end of the chapter, which will be mostly of 
  albegraic nature.  
}

\section{Basic mathematical notation}
\label{sec:notation}

This section is there to cover essential notation and define variables
for the use throughout the whole book. There will be some special 
abbreviations that need some time to get used to but will eventually 
make sense and lower the amount of writing in equations. In addition, 
there will be a list of all important functions and variables that will be 
defined properly.

\subsection{Einstein Summation Convention}
\label{subsec:summation-convention}

The summation convention proposed by Albert Einstein is helping researchers 
on the general theory of relativity to minimize their writing effort.
It is meant to remove all unnecessary summation terms and replace them solely 
with indices.

Let $\vec{a}$, $\vec{b}$ and $\vec{c}$ be vectors with four \emph{components}
and $\tens{M}$ a matrix with four rows and columns\footnote{
Notice the bold face notation for vectors and the sans serif notation 
for matrices.}. Then 
\begin{equation}
  \vec{a} = \vec{b} + \tens{M}\vec{c}
  \label{eq:matrix-equation}
\end{equation}
expresses this relationship in a \emph{matrix equation}.

In detail, take the matrix equation 
\begin{equation}
  \begin{pmatrix}
    a_1 \\ a_2 \\ a_3 \\ a_4
  \end{pmatrix}
  =
  \begin{pmatrix}
    b_1 \\ b_2 \\ b_3 \\ b_4
  \end{pmatrix}
  +
  \begin{pmatrix}
    M_{11} & M_{12} & M_{13} & M_{14} \\
    M_{21} & M_{22} & M_{23} & M_{24} \\
    M_{31} & M_{32} & M_{33} & M_{34} \\
    M_{41} & M_{42} & M_{43} & M_{44} \\
  \end{pmatrix}
  \begin{pmatrix}
    c_1 \\ c_2 \\ c_3 \\ c_4 
  \end{pmatrix}
  \label{eq:example-matrix-equation}
\end{equation}
as an example, which can be rewritten as 
\begin{equation}
  \begin{aligned}
    a_1 = b_1 + M_{11}c_1 + M_{12}c_2 + M_{13}c_3 + M_{14}c_4 \\
    a_2 = b_2 + M_{21}c_1 + M_{22}c_2 + M_{23}c_3 + M_{24}c_4 \\
    a_3 = b_3 + M_{31}c_1 + M_{32}c_2 + M_{33}c_3 + M_{34}c_4 \\
    a_4 = b_4 + M_{41}c_1 + M_{42}c_2 + M_{43}c_3 + M_{44}c_4 
  \end{aligned}
  \label{eq:rewritten-matrix-equation}
\end{equation}
and reduced to 
\begin{equation}
  a_i = b_i + \sum_{j=1}^{4} M_{ij}c_j.
  \label{eq:reduction-matrix-equation}
\end{equation}
By applying the \emph{Einstein Summation Convention} defined above and
leaving out the indices the equation~\ref{eq:matrix-equation} 
will be left~\cite{Article:summation_convention_01_barr}.

In proper writing, used throughout the rest of this book, it will 
reduce to 
\begin{equation}
  a_i = b_i + M_{ij}c_j 
  \label{eq:proper-summation-convention}
\end{equation}
with the use of the Einstein Summation 
Convention~\cite{Article:summation_convention_02_lei}.

\subsection{Derivatives}
\label{subsec:derivatives}

The fundamental concept of derivatives should be familiar but for 
the detailed notation used in this book a definition of some 
symbols is necessary.

\paragraph{Partial Derivative}

Usually, the partial derivative is expressed as 
\begin{equation}
  \frac{\partial}{\partial x}f(x,y)
  \label{eq:partial-derivative}
\end{equation}
for a function $f: \mathbb{R}^2 \rightarrow \mathbb{R}$. 
For the general theory of relativity this is expanded with the use 
of the coordinate basis\footnote{More on this in~\ref{sec:manifolds}} 
for an operation on any tensorial object, such as $T_{\mu\nu}$ 
for example. Therefore, the partial derivative shall herein and 
after be written as 
\begin{equation}
  \frac{\partial T_{\mu \nu}}{\partial x^{\mu}}.
  \label{eq:partial-derivative-gr}
\end{equation}

So as to eliminate the unnecessary effort for writing the whole expression 
over and over again, the partial derivative with respect to the 
coordinate basis may and will also be written as 
\begin{equation}
  \frac{\partial T_{\mu \nu}}{\partial x^{\mu}} = \partial_{\mu}T_{\mu \nu}
  \label{eq:partial-short}
\end{equation}
or, if preferred in the shortest version, as 
\begin{equation}
  \partial_{\mu}T_{\mu \nu} = T_{\mu \nu ,\mu}
  \label{eq:partial-very-short}
\end{equation}
with the use of a comma separating the indices from the index, which the term is 
derived with respect to~\cite{Book:general_relativity_01_wald,Book:general_relativity_02_dirac,Book:general_relativity_03_carroll}.

\paragraph{Covariant Derivative}

This new form of a derivative is needed for the correct differentiation 
in the general theory of relativity. Is is composed of the generic 
partial derivative and a correcting term. This correcting term will 
be explained in detail in later sections but for now the reader shall 
just keep in mind that it is meant to correct for deviations of the 
partial derivative in curved spacetimes.

Generally, this new term is expressed as 
\begin{equation}
  \nabla_{\mu}V^{\nu} = \partial_{\mu}V^{\nu} + \Gamma^{\nu}_{\;\;\mu \sigma}V^{\sigma}
  \label{eq:covariant-derivative}
\end{equation}
where $V^{\nu}$ are the components of a vector and 
$\Gamma^{\nu}_{\mu \sigma}$ denotes the correcting term.

In some books\footnote{Specifically, in~\cite{}}, this covariant derivative 
is also expressed in a shortened form, which in this case is 
\begin{equation}
  \nabla_{\mu}V^{\nu} = V^{\nu}_{\;\;;\mu}.
  \label{eq:covariant-very-short}
\end{equation}

\subsection{Operators}
\label{subsec:operators}

Some other preliminaries are necessary in order to provide a basic foundation 
of the notation used in this book. There are quite some definitions to be 
covered but with partial and covariant derivative out of the way and defined
in Subsection~\ref{subsec:derivatives} there are not as many left.

\paragraph{Laplacian Operator}



\subsection{Definitions}
\label{subsec:definitions}

\section{Transformations}
\label{sec:transformations}

Transformations are, out of necessity, covered in this section and there 
will be graphics provided for the better understanding of the topic.
Transformations are vital in the general concept of relativity because 
nothing is absolute anymore, so transformations will eventually take over 
as the standard method of changing the point of view. In later chapters,
the reader may go back to this part and check his understanding, 
when e.g. the cone of causality is introduced.

\subsection{Translations}
\label{subsec:translations}

\subsection{Rotations}
\label{subsec:rotations}

\subsection{Galilean Transformation}
\label{subsec:galilean}

\subsection{Lorentz Transformation}
\label{subsec:lorentz}

\section{Linear Algebra}
\label{sec:linear-algebra}

Although the basic understanding of Linear Algebra is a requirement, some
concepts will be explained in this section as well. These concepts include 
the vector space, basis vectors, matrix inversions and, subsequently, the
Gaussian Elimination algorithm. Furthermore, the concept of coordinate 
systems in both $\mathbb{R}^2$ and $\mathbb{R}^3$ will be covered.

\subsection{The vector space}
\label{subsec:vector-space}

\subsection{Matrices}
\label{subsec:matrices}

\subsection{Gaussian Elimination}
\label{subsec:gaussian-elimination}

\subsection{Cartesian Coordinates}
\label{subsec:cartesian}

\subsection{Spherical Coordinates}
\label{subsec:spherical-coordinates}

\section{Tensor Algebra}
\label{sec:tensor-algebra}

The foundation of the General Theory of Relativity is tensor algebra.
For the purpose of this book, basic concepts, such as vectors, dual vectors, tensors and the tensor product will be introduced and explained for further 
use in later chapters.

\subsection{Vectors}
\label{subsec:vectors}

\subsection{Dual Vectors}
\label{subsec:dual-vectors}

\subsection{Tensors}
\label{subsec:tensors}

\subsection{Tensors operations}
\label{subsec:tensor-operations}

\section{Analysis}
\label{sec:analysis}

Analysis is a prerequisite for reading this book. However, there are still
some topics to be covered in this book, which are Lie derivatives and 
Jacobian matrices. 

\subsection{Lie derivatives}
\label{subsec:lie}

\subsection{Jacobian}
\label{subsec:jacobian}

\section{Integration Methods}
\label{sec:integration}

Integration methods are basic methods for using analytic concepts in a 
computer graphics context. For the purpose of this book, four of these 
methods will be elaborated.

\subsection{Euler Integration}
\label{subsec:euler}

\subsection{Runge Kutta fourth order}
\label{subsec:runge-kutta}

\subsection{Runge Kutta Fehlberg fourth order}
\label{subsec:runge-kutta-fehlberg}

\subsection{SODA}
\label{subsec:soda}

%%%%%%%%%%%%%%%%%%%%%%%%%%%%%%%%%%%%%%%%%%%%%%%%%%%%%%%%%%%%%%%%%%%%%%%%%%%
% References
%%%%%%%%%%%%%%%%%%%%%%%%%%%%%%%%%%%%%%%%%%%%%%%%%%%%%%%%%%%%%%%%%%%%%%%%%%%

\bibliographystyle{spmpsci}
\bibliography{theory/ref-math}
